\subsection{Intermediate Proposals and Categorical Zoning Deltas}

To comprehensively satisfy the causal assumption that spatial dose-response architectures must model intermediate outcomes free from look-ahead bias, a longitudinal Natural Language Processing (NLP) pipeline was engineered to extract time-varying zoning trajectories from archival hearing transcripts. 

\subsubsection{Chronological Extraction and Zero-Leakage Join}
A multi-stage regex state machine was executed across 15,375 raw City Council and Planning Commission transcripts to capture semantic requests and administrative recommendations (e.g., \textit{Requested Zoning}, \textit{Staff Recommendation}, and \textit{Approved Zoning}). To ensure temporal fidelity, these extractions were formally mapped to their exact chronological \textit{source date} using historical meeting agendas. 

The resulting extractions were then dynamically injected into the biweekly causal panel using a strict backward-filling temporal match. This architectural decision mathematically guarantees that the deep learning sequence models only evaluate zoning proposals and applicant compromises that were formally documented \textit{on or before} a given 14-day biweekly period, fully eliminating target leakage.

\subsubsection{Quantitative Translation via Land Development Code}
To convert qualitative text strings (e.g., "LO-MU") into continuous quantitative targets for the Neural Ordinary Differential Equations, an explicit mapping dictionary based on Austin's Land Development Code (Chapter 25-2) was implemented. At every biweekly period, semantic zoning requests are automatically translated into their maximum allowable height (in feet) and maximum Floor-to-Area Ratio (FAR). If a transcript explicitly documents a negotiated variance (e.g., "height reduced to 35 feet"), this explicit numerical limit overrides the baseline dictionary value. This continuous temporal mapping captures the incremental degradation of allowable density caused by neighborhood mobilization.

\subsubsection{Validating Compromise Incidence}
Table~\ref{tbl:zoning_deltas} demonstrates the empirical validity of this approach. By isolating the final period of the causal panel, we observe that cases receiving formal protest petitions are overwhelmingly forced into intermediate negotiation cycles that end in a compromise or denial (where the Final Zoning diverges from the Initial Request).

\begin{table}[h!]
\centering
\begin{tabular}{l c c}
\toprule
\textbf{Petition Status} & \textbf{Cases (N)} & \textbf{Final Zoning Diverged (Compromise/Denial)} \\
\midrule
Unprotested (0) & 4,420 & 51.8\% \\
Protested (1) & 41 & 97.5\% \\
\bottomrule
\end{tabular}
\caption[Zoning Trajectory Deltas by Protest Status]{Zoning Trajectory Deltas by Protest Status. Protested cases show a near-universal rate (97.5\%) of final zoning outcomes diverging from the applicant's initial request, substantiating the resistance effect of NIMBY mobilization.}
\label{tbl:zoning_deltas}
\end{table}
